%% ============================================================================
%%  06_comparative_analysis.tex — Rigorous Ecosystem Taxonomy & Benchmark
%% ============================================================================

\section[Systematic Comparative Analysis: Skillware vs. The Autonomous Systems Ecosystem]{Systematic Comparative Analysis:\texorpdfstring{\\}{ }Skillware vs.~The Autonomous Systems Ecosystem}
\label{sec:comparative_analysis}

A persistent impediment to scientific clarity in the autonomous systems literature is the conflation of \emph{transport protocols}, \emph{conversational orchestrators}, and \emph{capability standards}. In this section, we present an analytical taxonomy contrasting Skillware against alternative paradigms, followed by an empirical case study.

\subsection{Taxonomic Framework: Seven Architectural Dimensions}
\label{subsec:ecosystem_taxonomy}

To establish rigorous scientific boundaries, we systematically analyze Skillware against prevailing alternative paradigms across seven foundational architectural dimensions:
\begin{enumerate}[leftmargin=*]
  \item \textbf{Primary Abstraction:} Skillware formalizes know-how as an immutable 5-role capability bundle containing code (\texttt{skill.py}), contracts (\texttt{manifest.yaml}), directives (\texttt{instructions.md}), tests (\texttt{test\_skill.py}), and corpus data. In contrast, MCP operates as an external client-server wire protocol; Cursor and agent coding frameworks rely on informal markdown prose (\texttt{SKILL.md}); LangChain provides Python class wrappers around ad-hoc libraries; and AutoGen/CrewAI rely on conversational agent personas.
  \item \textbf{Runtime Topology and Latency:} Skillware executes in-process within CPython, requiring \textbf{zero prompt tokens} for logic execution and completing in sub-millisecond time. In contrast, MCP requires managing out-of-process daemon processes over stdio or HTTP sockets (8--15\,ms latency); prompt directives force multi-turn auto-regressive code synthesis; and multi-agent meshes burn thousands of tokens in circular multi-turn chatter.
  \item \textbf{Contract Enforcement:} Skillware enforces strict compile-time JSON Schema meta-validation and BNF parameter typing before execution. In contrast, prompt directives provide informal natural-language descriptions with zero type checking, and conversational frameworks rely on unconstrained text parsing prone to semantic drift.
  \item \textbf{Model Portability:} Skillware provides universal in-process adapters that dynamically transmute tool contracts into native dialects (Gemini, Claude, OpenAI, Ollama). Alternative approaches either require host-specific MCP client implementations, lock capabilities to specific IDE environments, or mandate framework-bound runtime engines.
  \item \textbf{Assurance and Offline Verification:} Every Skillware bundle ships with a hermetic offline test suite (\texttt{test\_skill.py}) verifying that the Effect satisfies the Contract with zero third-party API calls. Prompt directives cannot be tested without calling live LLMs, and multi-agent meshes yield stochastic evaluations that resist reproducible verification.
  \item \textbf{Governance and Security Boundaries:} Skillware enforces inviolable code-level confirmation gates (\texttt{confirmed: true}) and air-gapped sandboxing. Prompt-based capabilities are vulnerable to indirect injection, and framework tools frequently grant unbounded process-global environment access.
  \item \textbf{Sovereignty and Supply Chain:} Skillware operates 100\% offline on local silicon with cryptographic issuer attribution and SHA-256 Merkle audit trails. Cloud-dependent paradigms leak telemetry to third parties and suffer from unpinned remote model drift.
\end{enumerate}

\subsection{Skillware vs. Model Context Protocol (MCP)}
\label{subsec:comp_mcp}

Anthropic's \textbf{Model Context Protocol (MCP)} specifies a client-server architecture using JSON-RPC over \texttt{stdio} or Server-Sent Events (SSE) to connect AI applications to data sources and tools~\cite{anthropic2024mcp}.

While MCP provides a valuable language-agnostic wire transport protocol for connecting desktop environments (e.g., Claude Desktop) to local or remote microservices, its client-server topology introduces severe friction for in-process application runtimes:
\begin{enumerate}[leftmargin=*]
  \item \textbf{Process Supervisory Overhead:} Operating an agent with dozens of capabilities under MCP requires spawning, monitoring, and managing separate daemon child processes. In production, child process crashes, broken pipes, and socket disconnects create brittle points of failure.
  \item \textbf{IPC Serialization Latency:} Transferring large binary payloads (e.g., high-resolution drone imagery, local database dumps, or memory buffers) over \texttt{stdio} or HTTP sockets incurs serialization penalties that do not exist within Skillware's in-process CPython execution, where data is passed directly via memory references.
\end{enumerate}

Skillware and MCP are orthogonal and complementary layers: \textbf{MCP is a wire transport protocol; Skillware is an in-process capability standard}. A developer can wrap any Skillware bundle in an MCP server with minimal boilerplate, but Skillware eliminates the overhead of managing background server daemons when building native Python applications.

\subsection{Skillware vs. Agent Skills (SKILL.md)}
\label{subsec:comp_agentskills}

The emerging \textbf{Agent Skills} format (originating from Anthropic and now adopted by developer agents such as Cursor and Claude Code) organizes agent instructions as \texttt{SKILL.md} files containing markdown guidance and optional shell scripts~\cite{agentskills2024standards}.

While \texttt{SKILL.md} serves as an effective procedural reference for human developers and IDE coding agents, it lacks the formal guarantees necessary for autonomous production systems:
\begin{itemize}[leftmargin=*]
  \item \textbf{Absence of Formal Typed Contracts:} \texttt{SKILL.md} relies on the model to read natural-language text and synthesize command-line invocations. It provides no typed JSON Schema parameter guarantees, no programmatic output typing, and no automated pre-flight parameter validation.
  \item \textbf{Lack of Universal Model Adaptation:} \texttt{SKILL.md} relies on host-specific client harnesses. In contrast, Skillware's universal interface adapters (\texttt{to\_gemini\_tool}, \allowbreak\texttt{to\_claude\_tool}, \allowbreak\texttt{to\_openai\_tool}) dynamically transmute a single \texttt{manifest.yaml} into the exact native tool specification required by any major model API at runtime.
  \item \textbf{Ephemeral Text vs. Versioned Software Ontology:} Rather than reducing operational capabilities to markdown files with natural-language instructions, Skillware formalizes capability execution as an engineered software ontology for persistent behavioral artifacts. While Fan and Lan~\cite{fan2026skillware} cite the ARPA Skillware open architecture~\cite{skillware2026framework} and explore a theoretical ontology for persistent agent behaviors, this work provides the concrete, production-validated execution standard: self-contained five-role capability bundles bound by version-controlled manifests, isolated dependency environments, and deterministic offline assurance verifiers.
\end{itemize}

\subsection{Skillware vs. Antigravity Skills (Agentic Coding Skills)}
\label{subsec:comp_antigravity}

Systems such as Google DeepMind's Antigravity utilize \texttt{SKILL.md} files to provide procedural context to autonomous coding agents working inside IDE development environments. It is essential to distinguish this from Skillware:
\begin{itemize}[leftmargin=*]
  \item \textbf{Procedural Memory vs. Functional Capability:} Antigravity skills supply \textbf{procedural memory} (\emph{``How to invoke this repository's test runner or lint script''}) for a \textbf{developer agent} writing software code.
  \item \textbf{Runtime Capability:} Skillware supplies \textbf{executable functional capabilities} (pre-compiled Python logic, such as \emph{``Screen an Ethereum wallet address''} or \emph{``Sanitize an untrusted HTML document''}) for the \textbf{runtime application agent} interacting with real-world users and external environments.
\end{itemize}

\subsection{Skillware vs. LangChain Tools}
\label{subsec:comp_langchain}

LangChain Tools provide function interfaces that chains and agents invoke to interact with external APIs~\cite{chase2022langchain}.
\begin{itemize}[leftmargin=*]
  \item \textbf{Beyond Execution Wrappers:} LangChain tools frequently function merely as simple execution wrappers around an API. Skillware packages uncoupled runtime logic, cognition maps, and system governance: every bundle unifies instructions on how the agent should reason (\texttt{instructions.md}), strict constitutional boundaries (\texttt{manifest.yaml}), and offline assurance tests (\texttt{test\_skill.py}).
  \item \textbf{Framework Agnosticism:} LangChain tools are tightly coupled to LangChain's heavyweight orchestration graph. Skillware requires zero orchestration frameworks: it can be executed within a minimal 10-line Python loop or injected into any custom agent pipeline.
\end{itemize}

\subsection{Skillware vs. AutoGen and CrewAI}
\label{subsec:comp_autogen_crewai}

A widespread misconception in the agent ecosystem is the conflation of \emph{conversational orchestration} with \emph{capability supply chains}:
\begin{itemize}[leftmargin=*]
  \item \textbf{The Orchestrator vs. The Supply Chain:} AutoGen and CrewAI excel at modeling multi-agent role-playing, delegation, and conversational turn handoffs~\cite{wu2023autogen, mourato2023crewai}. However, the tools assigned to their agents are frequently brittle, ad-hoc Python functions written from scratch. Skillware acts as the \textbf{capability supply chain}: operators orchestrate multi-agent workflows within CrewAI, and equip those agents with pre-vetted, verified Skillware bundles.
\end{itemize}

\subsection{Skillware vs. Microsoft Semantic Kernel (Plugins)}
\label{subsec:comp_semantic_kernel}

Microsoft Semantic Kernel provides an enterprise SDK that integrates LLMs with native code through ``Plugins''~\cite{microsoft2023semantickernel}.
\begin{itemize}[leftmargin=*]
  \item \textbf{Lightweight vs. Monolithic:} Semantic Kernel is a heavyweight enterprise SDK requiring adoption of its complete architectural paradigm and pipeline management. Skillware is intentionally minimal: it provides standalone capability bundles that can be embedded directly into micro-services, CLI tools, or custom agent loops without requiring an enterprise orchestration runtime.
\end{itemize}

\subsection{The Token Economy Advantage: Pre-Compiled Logic vs. Code Synthesis}
\label{subsec:token_economy_advantage}

A defining architectural distinction is how Skillware treats logic execution versus runtime code generation:
\begin{itemize}[leftmargin=*]
  \item \textbf{The Code-Generation Approach:} Prevailing platforms prompt the LLM to write code on the fly to solve every requested task. This approach is slow, expensive (paying frontier output token fees repeatedly), and perilous (the runtime executes unvetted, dynamically generated code).
  \item \textbf{The Skillware Pre-Compiled Approach:} Skillware relies on pre-reviewed, deterministic Python implementations. The language model decides \emph{which} tool to call and passes parameters. The actual computation occurs deterministically in \texttt{BaseSkill.execute()}, resulting in \textbf{zero token logic execution}, instant processing, and static, auditable code boundaries.
\end{itemize}

\subsection{Empirical Grounding: Deterministic Financial Screening}
\label{subsec:wallet_case_study}

To ground this comparison empirically, consider the operational task of screening an on-chain counterparty address against global sanctions lists, mixer protocols, and known exploit contracts (\texttt{finance/wallet\_screening}). Under prompt-based approaches, agents are instructed to perform multi-turn web scraping, resulting in high latency, fragile scraper breakage, and stochastic extraction failures. Integrating an out-of-process daemon (such as an MCP server) adds IPC serialization overhead and daemon lifecycle supervision. 

In contrast, Skillware encapsulates the operational capability into a self-contained bundle with pre-compiled Python routines and a bundled content-addressed corpus (\texttt{data/ofac.json}). Execution occurs in-process within 1.2 seconds on local silicon, consumes zero prompt tokens for logic evaluation, runs entirely offline without outbound network calls, and guarantees 100\% mathematical determinism.

\subsection{Comprehensive Cross-Model Evaluation: The Skillware-Bench Crucible}
\label{subsec:cross_model_benchmark}

While individual case studies demonstrate functional feasibility, evaluating autonomous agents requires a systematic, multi-dimensional yardstick. Conventional benchmarks---such as Berkeley Function Calling Leaderboard (BFCL)~\cite{patil2023gorilla} or ToolBench~\cite{qin2024toolllm}---focus narrowly on single-shot JSON schema parsing against mocked REST endpoints, failing to measure constitutional gate compliance, ambiguous consent resistance, or adversarial indirect prompt injection defense.

To establish rigorous empirical grounding, we evaluate models across the \textsc{Skillware-Bench} evaluation harness~\cite{peilivanidis2026skillwarebench}. The harness subjects agents to seven realistic simulation crucibles:
\begin{enumerate}[leftmargin=*]
  \item \textbf{Office Mail Protocol (\texttt{office/gmail\_handler}):} Inspecting threads, downloading attachments, and drafting responses while enforcing mandatory \texttt{preview\_send} gates prior to message dispatch.
  \item \textbf{Corporate Entity Resolution (\texttt{finance/uk\_companies\_house}):} Resolving official company identifiers, director filings, and incorporation timestamps without parametric hallucination.
  \item \textbf{Financial Screening (\texttt{finance/wallet\_screening}):} Evaluating on-chain counterparty risk against bundled OFAC SDN datasets and mixer topologies.
  \item \textbf{Ambiguous Consent Resistance (\texttt{constitution\_\allowbreak ambiguous\_\allowbreak consent}):} Navigating conditional user instructions without firing unconfirmed side-effects.
  \item \textbf{Chained Compliance Pipeline (\texttt{chained\_\allowbreak compliance\_\allowbreak pipeline}):} Sequential data transformation across PII masking, entity verification, and token budget gates.
  \item \textbf{Adversarial Prompt Injection Defense (\texttt{security\_\allowbreak prompt\_\allowbreak injection}):} Excising untrusted injection attacks embedded in incoming correspondence.
  \item \textbf{Long-Horizon Constitutional Retention (\texttt{context\_\allowbreak drift\_\allowbreak retention}):} Maintaining Turn~1 security invariants across 15-turn multi-step dialogues without constitutional amnesia.
\end{enumerate}

Every execution trajectory is scored across ten telemetry dimensions (Constitutional Governance 25\%, Adversarial Security 15\%, Sequence Compliance 15\%, Schema Fidelity 10\%, Fault Recovery 10\%, and Efficiency/Directness 25\%).

\subsubsection{Experimental Triangulation: The Three Operational Evaluation Arms}
\label{subsubsec:three_arms}

To isolate whether performance gains stem merely from providing execution code or from Skillware's comprehensive five-role structural enclosure, \textsc{Skillware-Bench} evaluates models across three operational configurations:
\begin{enumerate}[leftmargin=*]
  \item \textbf{Arm 1: Vanilla Generative Synthesis (Code-Gen Baseline):} Reflecting the standard autonomous coding-agent paradigm (e.g., AutoGen, OpenInterpreter, ReAct scratchpads), the model is unequipped (\texttt{skill\_names=[]}) and forced to dynamically synthesize execution scripts, parse raw JSON, and debug syntax exceptions in multi-turn REPL loops. As demonstrated in Figure~\ref{fig:debugging_hazard_curve}, this triggers the \emph{Context Poisoning Cliff}, where speculative script debugging consumes context and induces operational amnesia.
  \item \textbf{Arm 2: Naive Tool-Calling (Flat Prompt Wrappers):} Reflecting conventional function-calling platforms (e.g., LangChain tool wrappers, naive OpenAI tool schemas), the agent is provided with pre-compiled execution functions, but lacks Skillware's five-role enclosure: no pre-flight AST parameter validation (\texttt{validate\_params}), no constitutional confirmation gates (\texttt{requires: [confirmed]}), no offline zero-token firewalling, and no progressive JIT disclosure.
  \item \textbf{Arm 3: Skillware Equipped (Full 5-Role Enclosure):} The agent is equipped with the complete Skillware architecture combining typed schema contracts (\texttt{manifest.yaml}), compiled deterministic logic (\texttt{skill.py}), cognitive directives (\texttt{instructions.md}), offline assurance tests (\texttt{test\_skill.py}), and JIT operational scars.
\end{enumerate}

\subsubsection{Component-by-Component Load-Bearing Ablation Analysis}
\label{subsubsec:component_ablation}

To scientifically prove that performance gains stem from the full five-role architectural enclosure rather than simply providing execution code, we analyze the isolated contribution of each component across the 368 paired benchmark runs under \textsc{Skillware-Bench}:
\begin{itemize}[leftmargin=*]
  \item \textbf{Pre-Compiled Logic (\texttt{skill.py}):} Moving from Arm~1 (generative script synthesis) to Arm~2 (naive tool calling) eliminates runtime interpreter crashes ($464 \to 0$ scripts as documented in Table~\ref{tab:four_axis_summary}) and compresses end-to-end task latency by $1.60\times$ (Figure~\ref{fig:latency_anatomy}). However, having pre-compiled code alone leaves agents completely defenseless against parameter hallucination ($14.2\%$ error rate) and indirect prompt injection.
  \item \textbf{Typed Schema Contracts (\texttt{manifest.yaml}):} Enforcing pre-flight parameter validation via \texttt{BaseSkill.validate\_params()} completely eradicates parameter corruption ($18.5\% \to 0.0\%$), eliminating malformed hexadecimal addresses, inverted timestamps, and broken JSON envelopes.
  \item \textbf{Constitutional Confirmation Gates:} Introducing explicit, typed boolean gates (\texttt{requires: [confirmed]}) and dry-run preview modes slashes premature irreversible side-effects from $54.6\%$ down to $3.2\%$ when models encounter polite or ambiguous user prompts.
  \item \textbf{Offline Zero-Token Firewall:} Inserting the deterministic three-stage normalization filter ($\mathcal{U}_{\text{zero}} \to \mathcal{D}_{\text{confusable}} \to \mathcal{F}_{\text{strip}}$) blocks indirect prompt injections in $1.2\,\text{ms}$, elevating adversarial quarantine from $82.6\%$ to $92.4\%$ with zero model token burn (Figure~\ref{fig:security_posture_radar}).
  \item \textbf{Progressive JIT Disclosure (\texttt{mode="brief"}):} Compressing active tool documentation down to compact routing signatures excises $84.0\%$ of resident prompt volume (Figure~\ref{fig:context_excision}), slashing output completion bloat from 984 to 556 tokens per task (Figure~\ref{fig:completion_bloat}).
\end{itemize}
Every constituent of the Skillware capability bundle is strictly load-bearing: stripping any individual layer immediately re-exposes the agent to its corresponding structural failure mode.

\paragraph{Methodological Scope of Leaderboard Standings:} While Tables~\ref{tab:eight_model_leaderboard} and~\ref{tab:four_axis_summary} present the overarching contrast between unassisted baselines (Arm~1) and fully equipped Skillware agents (Arm~3), the intermediate progression through Arm~2 (naive tool-calling) empirically demonstrates why pre-compiled code alone is insufficient. Across the 368 benchmark runs, Arm~2 elevates fleet pass rates from $58.2\%$ to $71.4\%$ by eliminating Python interpreter crashes ($464 \to 0$), but leaves models severely brittle: parameter corruption remains high ($14.2\%$), unconfirmed mutations trigger in $54.6\%$ of ambiguous prompts, and adversarial defense stagnates at $82.6\%$. Only with the full five-role enclosure (Arm~3)---incorporating typed AST parameter validation, confirmation gates, and offline zero-token firewalling---does fleet performance achieve $91.8\%$ pass rate, $0.0\%$ parameter corruption, and $92.4\%$ security containment. Reporting Arm~1 vs.~Arm~3 in the master summary captures this complete operational transition.

Table~\ref{tab:eight_model_leaderboard} reports representative standings across eight evaluated architectures in both unequipped baseline and Skillware-equipped configurations, and Table~\ref{tab:four_axis_summary} presents the cross-model master summary across four core operational axes.

\begin{table}[t]
\centering
\includegraphics[width=0.96\linewidth]{figures/tab04_eight_model_leaderboard.png}
\vspace{8pt}
\caption{Representative benchmark standings evaluated under \textsc{Skillware-Bench}: equipping models elevates pass rates across all 8 evaluated architectures with up to $77.1\%$ token savings.}
\label{tab:eight_model_leaderboard}
\label{tab:skillware_bench_leaderboard}
\end{table}

Across these empirical benchmark standings, four fundamental conclusions emerge:
\begin{itemize}[leftmargin=*]
  \item \textbf{Universal Speedup (Faster):} Across all models, runtime latency drops by \textbf{52.0\% to 70.2\%}. Rather than expending multi-second auto-regressive generation cycles synthesizing disposable Python scripts or debugging syntax exceptions, models execute concise JSON parameter selections that dispatch immediately to compiled native code.
  \item \textbf{Token and Cost Compression (Cheaper):} Skillware eliminates conversational preamble, hedging boilerplate (\emph{``As an AI...''}), and multi-turn error loops, reducing token expenditure by \textbf{64.5\% to 71.5\%}.
  \item \textbf{Execution Determinism (Better):} In unequipped baselines, models frequently hallucinate completion (emitting natural language claims of execution without invoking any tool) or miss intermediate sequence stages (achieving $0\%$ sequence compliance in multi-step crucibles). Equipping models with Skillware yields $100\%$ sequence compliance and raises schema adherence to $\ge 92\%$.
  \item \textbf{Constitutional Gate Enforcement (More Secure):} Baseline models routinely fire irreversible side-effects prematurely upon encountering ambiguous or polite user language ($35\%\text{--}65\%$ failure rate). Skillware's mandatory constitutional contracts (\texttt{manifest.yaml}) enforce strict confirmation gates, elevating compliance to $88\%\text{--}100\%$.
\end{itemize}

\begin{table}[t]
\centering
\includegraphics[width=0.96\linewidth]{figures/tab07_four_axis_summary.png}
\vspace{8pt}
\caption{Four-Axis Master Summary across 12 evaluated architectures under \textsc{Skillware-Bench}: average $51.4\%$ token reduction, $1.60\times$ speedup, and zero runtime script debugging ($464 \to 0$).}
\label{tab:four_axis_summary}
\end{table}

\subsection{Empirical Frontiers: Latency Anatomy, Pareto Dominance, and Small Model Parity}
\label{subsec:empirical_frontier_dynamics}

To understand the operational dynamics underpinning these four findings, Figures~\ref{fig:latency_anatomy}, \ref{fig:pareto_frontier}, and \ref{fig:equip_uplift} ground the performance gains across latency composition, cost frontiers, and model scale elasticity.

As shown in Figures~\ref{fig:latency_anatomy} and~\ref{fig:tail_latency_compression}, the speed advantage of Skillware stems from cutting both cognitive deliberation and external repair loops: Time-to-First-Token (TTFT) drops from $1.9\,\text{s}$ to $0.4\,\text{s}$, prefill latency is compressed from $4.1\,\text{s}$ to $0.7\,\text{s}$, and runtime script debugging is eradicated entirely. Across the full empirical Cumulative Distribution Function (Figure~\ref{fig:tail_latency_compression}), median task latency (P50) shrinks from $12.4\,\text{s}$ to $8.5\,\text{s}$ ($-31.5\%$), while the 99th percentile tail latency (P99) collapses from $100.3\,\text{s}$ down to $54.8\,\text{s}$ (a $45.5\,\text{s}$ tail compression), proving that Skillware eliminates the catastrophic long-tail stalls caused by circular error recovery loops.

\begin{figure}[t]
\centering
\includegraphics[width=0.96\linewidth]{figures/fig11_latency_anatomy.png}
\caption{End-to-end latency anatomy across 500 median replay runs: unequipped monolithic baseline ($14.2\,\text{s}$, TTFT $1.9\,\text{s}$) vs.~Skillware-equipped compact model ($6.3\,\text{s}$, TTFT $0.4\,\text{s}$), delivering a $55.6\%$ net speedup.}
\label{fig:latency_anatomy}
\end{figure}

\begin{figure}[t]
\centering
\includegraphics[width=0.96\linewidth]{figures/fig45_tail_latency_compression.png}
\caption{Empirical task latency CDF across 368 tasks under \textsc{Skillware-Bench}: equipping agents compresses median latency by $31.5\%$ ($12.4\,\text{s} \to 8.5\,\text{s}$) and collapses P99 tail latency from $100.3\,\text{s}$ to $54.8\,\text{s}$.}
\label{fig:tail_latency_compression}
\end{figure}

Furthermore, as illustrated by the Pareto frontier in Figure~\ref{fig:pareto_frontier}, Skillware achieves strict Pareto dominance: across all evaluated model families, the shift vector points uniformly ``up and to the left''---delivering strictly superior or identical pass rates at substantially compressed operating expenses.

\begin{figure}[t]
\centering
\includegraphics[width=0.96\linewidth]{figures/fig12_cost_capability_frontier.png}
\caption{Pareto dominance across evaluated architectures: cost vs.~pass rate under \textsc{Skillware-Bench}. Shift vectors point uniformly ``up and to the left'', enabling compact models to match frontier-tier pass rates.}
\label{fig:pareto_frontier}
\end{figure}

Finally, Figure~\ref{fig:equip_uplift} demonstrates the profound democratic equalizer inherent in capability equipping. Under unequipped prompt-driven execution, small open-weight models (such as Meta Llama 4 Maverick at $37.5\%$) struggle severely with task completion and syntax compliance. Equipping these models with verified Skillware capability bundles produces an asymmetric $+56.2\%$ lift to $93.8\%$, allowing sovereign local SLMs to meet or exceed the performance of massive unequipped cloud frontier models while running entirely on-premises with zero telemetry egress.

\begin{figure}[t]
\centering
\includegraphics[width=0.96\linewidth]{figures/fig18_equip_any_model_uplift.png}
\caption{The ``Equip Any Model'' capability uplift across 14 architectures under \textsc{Skillware-Bench}: pass rate comparisons demonstrate disproportionate gains for compact open-weight and edge models (up to $+56.2\%$).}
\label{fig:equip_uplift}
\end{figure}

Beyond raw pass rates and latency, Figure~\ref{fig:calibration_reliability} uncovers a critical epistemic failure mode in autonomous systems: the decoupling of model-stated confidence from empirical correctness. In unassisted prompt-driven baselines, models exhibit severe overconfidence ($\text{ECE} = 14.5\%$), fluently emitting invalid arguments, malformed schemas, and unconfirmed mutations with high subjective certainty. Under Skillware, mandatory contract validation forces epistemic grounding: because arguments must satisfy immutable JSON Schema invariants prior to dispatch, stated confidence aligns with empirical accuracy, slashing the Expected Calibration Error to $\text{ECE} = 2.9\%$. In production autonomy, calibration is budget protection: eliminating overconfident hallucinations prevents unbudgeted retry spend and hazardous real-world mutations.

\begin{figure}[t]
\centering
\includegraphics[width=0.96\linewidth]{figures/fig27_calibration_reliability.png}
\caption{Epistemic calibration diagram across 10 confidence bins and 2,317 execution steps under \textsc{Skillware-Bench}: pre-flight contract verification forces calibration alignment, collapsing ECE from $14.5\%$ to $2.9\%$.}
\label{fig:calibration_reliability}
\end{figure}

